%! Algebra 2 — Unit 1, Lesson 3 — Exit Ticket
\documentclass[10pt]{article}
\usepackage{algebra2-article}
\usepackage{algebra2-boxes}

%=============================================================================
\begin{document}

\pageheader{Unit 1, Lesson 3}{Exit Ticket}
\namedateperiod

\vspace{0.3in}

\begin{enumerate}[leftmargin=*, itemsep=1.8em]

    % ── Question 1 ──────────────────────────────────────────────────────────
    \item For $f(x) = x^2 - 2x - 3$, identify each feature. \textbf{Show your work.}

    \vspace{0.1in}
    \writelines{4}

    \vspace{0.05in}
    Zeros: \blank{3cm} \quad Vertex: \blank{3cm} \quad $y$-intercept: \blank{3cm}

    % ── Question 2 ──────────────────────────────────────────────────────────
    \item A linear model is $y = 4.1x + 22$, where $x$ = months since January and
    $y$ = number of new customers.
    \textbf{Interpret the slope and the $y$-intercept in context.}

    \vspace{0.15in}
    \writelines{4}

    % ── Question 3 (SOL-style MC) ────────────────────────────────────────────
    \item Which description best fits a scatterplot that would be modeled by a \emph{linear} function?

    \vspace{0.2in}
    \begin{tabularx}{\linewidth}{X X}
        \textbf{A.}\quad Points curve upward, then level off
        & \textbf{B.}\quad Points are clustered near a straight diagonal line \\[12pt]
        \textbf{C.}\quad Points form a symmetric U-shape
        & \textbf{D.}\quad Points are randomly scattered with no visible pattern
    \end{tabularx}

\end{enumerate}

\end{document}
